% !TEX root = ../../tree_notes.tex
% =====================================================================
%  9.4  The Cartesian Tree as a Random Treap     sec:c26-treapview
%  Exports: prop:c26-carters-treap  thm:c26-randbst  fig:c26-iso
%           ex:c26-treapview-run
%           prob:c26-mbt2  lst:c26-mbt2-soln
% =====================================================================
\section{The Cartesian Tree as a Random Treap}\label{sec:c26-treapview}

\begin{previouslybox}
\begin{enumerate}[nosep, leftmargin=1.6em]
  \item A treap is a BST on \emph{keys} and a heap on \emph{priorities}. Given a fixed key set and a
        fixed priority per key, how many treaps realise them? (\chref{c06}, \defref{c06-treap} -- exactly
        one, when priorities are distinct)
  \item What does \thmref{c06-exptheight} guarantee for a treap with i.i.d.\ random priorities, and over
        what is the expectation taken? ($\bigO{\log n}$ height, over the random priorities)
  \item A Cartesian tree is a BST on which quantity and a heap on which? (\secref{c26-cartesian}: BST on
        \emph{index}, heap on \emph{value})
  \item Build a BST by inserting $3,1,4,2$ in that order. Is the root $3$ (the first key) or $1$ (the
        smallest)? (the first inserted -- $3$)
\end{enumerate}
\end{previouslybox}

\begin{topicintro}
This last section is a lens, not new machinery. Read a Cartesian tree one way and it is ``BST on index,
heap on value.'' Read the \emph{same} two-orders definition with the second order called a priority, and
it is \emph{exactly a treap} (\defref{c06-treap}) whose keys are the positions and whose priorities are
the array values. The tree of \secref{c26-cartesian} and the tree of Chapter~8 were never different
objects; only the story about what the heap-order ``means'' changed. Flipping the perspective once more
-- let the \emph{priority be the insertion time} instead of the value -- recovers a third classic: the
shape of a plain BST built by inserting a sequence is a Cartesian tree, so ``insert a random permutation
into a BST'' is ``build a Cartesian tree with random priorities,'' whose height is logarithmic \emph{in
expectation} by Chapter~8's theorem, quoted and never re-proved (\thmref{c26-randbst}).

We make the equivalence precise (\propref{c26-carters-treap}), show all three readings of one picture
(\figref{c26-iso}), and then cash it on LeetCode~998: append a value to the far right of a
maximum-tree's underlying array and re-root it -- which the treap view makes a one-line walk down the
right spine. \textbf{Prerequisites.} \secref{c26-cartesian} (the tree) and Chapter~8's treap definition
and height theorem (\chref{c06}) \emph{as black boxes}.
\textbf{After this section you can}
\begin{itemize}[nosep]
  \item state and prove that a Cartesian tree \emph{is} a treap with value-priorities, and vice-versa; \hfill\lv{3}--\lv{4}
  \item explain why inserting a random permutation yields a random treap and cite the expected height; \hfill\lv{4}
  \item solve right-end maximum-tree insertions (\probref{c26-mbt2}) from the priority viewpoint; \hfill\lv{4}
\end{itemize}
\end{topicintro}

\begin{recognitionbox}[Pattern card -- the same tree, wearing priorities]
\reckey{a structure that is ``ordered by one key and heap-ordered by a second number''; or ``insert these
into a BST, what shape results?''; or a \emph{max/min-tree} you must extend at one end.}
\recuse{
  \item "rebuild a maximum binary tree after appending a value" -- LC 998 (\probref{c26-mbt2}), the treap
        right-spine view;
  \item "expected depth / comparisons when inserting a random order into a BST" -- \thmref{c26-randbst},
        reduce to Chapter~8's bound;
  \item any treap you build by hand is secretly a Cartesian tree over the (key, priority) pairs.
}
\recnotuse{
  \item you need to \emph{maintain} balance under adversarial edits -- that is a real balanced tree with
        rotations (Chapter~8); priorities here are read off, not chosen to force balance;
  \item you want the treap \code{split}/\code{merge} engine -- that lives in \secref{c06-treaps}, this
        section borrows its \emph{definition} only;
  \item priorities must be \emph{random and independent} -- the array values here are given, not sampled.
}
One-line test: \emph{BST on one field, heap on another -- and the second field is really a priority? Then
you are looking at a treap in a Cartesian tree's clothes, and Chapter~8's guarantees apply for free.}
\end{recognitionbox}

\subsection{Intuition}

The definition \defref{c26-cartesian} demanded two things: inorder $=$ the index order, and a heap on the
values. The treap definition \defref{c06-treap} demands two things: a BST on the \emph{keys}, and a heap on
the \emph{priorities}. Set $\text{key} = \text{index}$ and $\text{priority} = \text{value}$ and the two
demands are word-for-word identical -- a BST on the index is precisely ``inorder reads the indices in
order,'' and a heap on the value is the Cartesian heap rule. Since Chapter~8 already proved a treap with
distinct priorities is \emph{unique}, that is a second proof of \thmref{c26-unique}: the Cartesian tree is
unique because it is a treap.

Now rotate the viewpoint. Suppose instead you do not care about values at all; you just \emph{insert keys
into a plain BST} in some arrival order. The element inserted first becomes the root; the ones after it
fall left or right by size, and recursively each subsequence rebuilds the same way. That is a Cartesian
tree whose \emph{priority is the insertion time} (smallest time $=$ highest in the tree). If the arrival
order is a uniformly random permutation, the insertion times are i.i.d.-in-effect random priorities -- so
the resulting BST has the \emph{same distribution as a random treap}, and everything Chapter~8 proved about
treap height transfers verbatim. That is \thmref{c26-randbst}, and its proof is a single sentence plus a
citation -- exactly the anti-repetition this chapter promised.

\begin{figure}[H]
\centering
\begin{tikzpicture}[diagram, scale=0.8, transform shape]
  % (1) min-Cartesian tree of A, heap on value
  \node[troot]  (c2) at (0,0)   {$1$};
  \node[tleaf]  (c1) at (-1.6,-1.2) {$3$};
  \node[tnode]  (c4) at (1.6,-1.2)  {$2$};
  \node[tleaf]  (c3) at (0.8,-2.4)   {$4$};
  \draw[tedgehl] (c2)--(c1); \draw[tedgehl] (c2)--(c4); \draw[tedgehl] (c4)--(c3);
  \node[annot, font=\scriptsize\bfseries] at (0,1.0) {(a) Cartesian: heap on value};
  % (2) treap: key=index (shown), priority=value (subscript)
  \node[troot]  (t2) at (6.5,0)  {$\mathit{2}$};
  \node[tleaf]  (t1) at (4.9,-1.2) {$\mathit{1}$};
  \node[tnode]  (t4) at (8.1,-1.2) {$\mathit{4}$};
  \node[tleaf]  (t3) at (7.3,-2.4) {$\mathit{3}$};
  \draw (t2)--(t1); \draw (t2)--(t4); \draw (t4)--(t3);
  \node[annot, font=\scriptsize] at (4.9,-1.2){\phantom{x}};
  \node[font=\scriptsize, text=darkgray] at (6.5,0.5) {\scriptsize p1};
  \node[annot, font=\scriptsize\bfseries] at (6.5,1.0) {(b) Treap: key=index, prio=value};
  \node[annot, align=left, font=\scriptsize] at (10.3,-1.2)
    {nodes $1..4$ by key;\\ priorities $a_2{=}1<a_4{=}2<a_1{=}3$\\ $\Rightarrow$ root $=$ key $2$};
  % (3) insertion-order BST (priority=index), array order 3,1,4,2 -> keys are VALUES, heap on index
  \node[troot]  (i3) at (0,-5.0) {$3$};
  \node[tnode]  (i1) at (-1.6,-6.2) {$1$};
  \node[tleaf]  (i4) at (1.6,-6.2) {$4$};
  \node[tleaf]  (i2) at (-0.8,-7.4) {$2$};
  \draw (i3)--(i1); \draw (i3)--(i4); \draw (i1)--(i2);
  \node[annot, font=\scriptsize\bfseries] at (0,-4.0) {(c) BST by insertion order $3,1,4,2$};
  \node[annot, align=left, font=\scriptsize] at (3.2,-6.2)
    {root $=$ first in $=$ value $3$\\ (min index). BST on value,\\ heap on insertion time $=$ index:\\ \emph{a different} Cartesian tree\\ of the same array};
\end{tikzpicture}
\caption{One array $A=(3,1,4,2)$, three readings. (a) min-Cartesian tree -- inorder $=$ index, heap on
value (\defref{c26-cartesian}). (b) the \emph{identical} tree as a treap: keys are the indices $1..4$
(labels here), priorities are the values, root $=$ least-priority key $=2$ (\defref{c06-treap});
(a)$\equiv$(b) is \propref{c26-carters-treap}. (c) is the \emph{other} assignment -- a plain BST built by
inserting $3,1,4,2$, which is a Cartesian tree keyed by value with priority $=$ index, the object that
makes a random permutation a random treap (\thmref{c26-randbst}).}\label{fig:c26-iso}
\end{figure}

\subsection{The equivalence}

\begin{proposition}{Cartesian tree $=$ treap; insertion BST $=$ Cartesian by time}{c26-carters-treap}
Let $a[1..n]$ have distinct values.
\begin{itemize}[nosep]
  \item[(i)] The min-Cartesian tree of $a$ is exactly the treap whose keys are $\{1,\dots,n\}$ and whose
        priority of key $i$ is $a[i]$. Conversely that treap is the Cartesian tree.
  \item[(ii)] Build a BST by inserting the values $a[1],\dots,a[n]$ in that \emph{arrival} order. The
        resulting tree is the Cartesian tree keyed by \emph{value} with priority equal to the insertion
        index $i$ (min-heap on $i$, since the earliest-inserted key sits highest).
\end{itemize}
\end{proposition}
\begin{proofbox}
(i) A treap is a BST on keys and a heap on priorities (\defref{c06-treap}). With key $=i$ and priority
$=a[i]$, "BST on keys" says the inorder lists keys ascending -- i.e.\ $1,2,\dots,n$, the Cartesian
inorder condition; "min-heap on priorities" says a node's priority $a[i]$ is $\le$ its children's, the
Cartesian heap condition. So the treap satisfies \defref{c26-cartesian} and, both being unique under
distinct values/priorities, they are the same tree. (ii) The first key inserted, $a[1]$, is the root (no
key preceded it). Recursively, within any subtree the \emph{earliest} member of that key-block to arrive is
its local root, and arrivals split left/right purely by comparison -- so the shape is a Cartesian tree on
the key set where the priority is arrival time (index), and the root is the minimum-index member: a
min-heap on index. Distinct priorities (indices) give uniqueness as in (i). \end{proofbox}

\noindent\textbf{In plain words.} Swap the roles of the two orders and the same picture is a Cartesian
tree, a value-priority treap, or an insertion-time BST. Chapter~8's treap uniqueness \emph{is}
\thmref{c26-unique}.

\begin{theorem}{A random permutation builds a random BST of logarithmic height}{c26-randbst}
Insert a uniformly random permutation of $n$ distinct keys into an initially empty binary search tree.
By \propref{c26-carters-treap}(ii) the resulting tree is a Cartesian tree whose priorities are the
permutation's arrival positions -- i.e.\ a treap with i.i.d.-in-effect distinct random priorities. Its
expected height is therefore $\bigO{\log n}$, and its expected search cost is $\bigO{\log n}$; both are
the direct content of Chapter~8's \thmref{c06-exptheight}, invoked as a black box and \emph{not}
re-derived here.
\end{theorem}
\begin{proofbox}
A uniform random permutation assigns to each key an arrival time that is a uniformly random injection
into $\{1,\dots,n\}$ -- the same as sampling distinct priorities in random order. By
\propref{c26-carters-treap}(ii) the insertion-order BST is the treap on those (key, priority) pairs, so its
distribution equals that of a random treap. \thmref{c06-exptheight} states a random treap has expected
height $c\,\ln n$ for a constant $c$ (and, more strongly, is \emph{balanced} with high probability). Since
a BST's every operation costs $\bigO{h}$ (\secref{c05-insert-delete}), expected cost is $\bigO{\log n}$. No
new proof is needed -- the reduction is the whole argument, which is why this lives in a section rather
than a chapter. \end{proofbox}

\begin{insightbox}
Three chapters now share one sentence. Chapter~8 chose priorities \emph{at random} to force a short tree
you can edit. This chapter's array \emph{hands you} priorities (the values), and the tree is whatever
those dictate -- which you can still exploit (histogram, RMQ) because a monotonic stack reads the same
priorities faster than any rotation. The unifying idea: \textbf{a BST-plus-heap on a second field is one
structure, and what you call the second field -- value, priority, or arrival time -- only picks the story.}
\end{insightbox}

\subsection{Worked Example}

\begin{example}{Reading $A$ as a treap and as an insertion BST}{c26-treapview-run}
For $A=(3,1,4,2)$: \figref{c26-iso}(a) and (b) are the same tree. As a \emph{treap} (b): keys $1,2,3,4$
with priorities $3,1,4,2$. The least priority is key $2$ (priority $1$), so it is the root; keys less
than $2$ (just key $1$) go left, keys greater ($3,4$) go right, and among them key $4$ (priority $2$) is
rooted above key $3$ (priority $4$) -- identical to the min-Cartesian build, as \propref{c26-carters-treap}
promises. As an \emph{insertion BST} (c): insert $3$, then $1$ (goes left of $3$), then $4$ (right of
$3$), then $2$ (left of $3$, right of $1$). Root $3$ because it arrived first (smallest index). Now the
keys are the \emph{values} and the priority is the \emph{index}: root $3$ has index $1$, its subtree keys
$\{1,2\}$ have larger indices, and $1$ (index $2$) precedes $2$ (index $4$), so $1$ is $2$'s parent. Same
rule, relabelled -- \propref{c26-carters-treap}(ii) on the nose. \checkmark
\end{example}

\subsection{Common Pitfalls}\label{sec:c26-treapview-pitfalls}

\begin{warningbox}[Trap 1 -- re-proving the treap height bound]
\thmref{c26-randbst}'s height statement is \emph{imported} from \thmref{c06-exptheight}; the anti-repetition
contract forbids rebuilding it. The only new content here is the equivalence \propref{c26-carters-treap}.
If you find yourself writing a fresh "expected height of random BST" proof, you have left the chapter.
\end{warningbox}

\begin{warningbox}[Trap 2 -- confusing which field is the priority]
In reading (a)/(b) the \emph{value} is the priority and the \emph{index} is the key (min-Cartesian /
value-treap). In reading (c) the \emph{index (arrival time)} is the priority and the \emph{value} is the
key. Both are valid Cartesian trees but of \emph{different} shapes -- \figref{c26-iso}(a) and (c) are not
the same tree. State which field is the priority before arguing anything.
\end{warningbox}

\begin{warningbox}[Trap 3 -- choosing priorities to force balance]
A treap picks priorities \emph{randomly} to \emph{guarantee} short height against adversarial keys. A
Cartesian tree gets its priorities \emph{from the data} and makes no such promise: a sorted array gives a
value-treap of height $n$ (a chain) -- the very worst case the randomised treap avoids. Never claim a
read-off Cartesian tree is balanced; only the \emph{random-permutation} BST of \thmref{c26-randbst} is
balanced in expectation.
\end{warningbox}

\subsection{Problem Lab}

The anchor is LeetCode~998 -- given the root of a \emph{maximum tree} (a max-Cartesian tree) and a value
appended to the far right of its hidden array, return the new maximum tree. The priority view makes it
trivial: the new value is the \emph{largest key} (last position), so it descends only along right
children, becoming a treap insert restricted to the right spine.

% ---------------------------------------------------------------------------
\solvedproblem{c26-mbt2}{Maximum Binary Tree II}{LeetCode 998}
  {https://leetcode.com/problems/maximum-binary-tree-ii/}{LC Medium · L3}
  {max-Cartesian / treap view, right-spine descent, priority as value}

\begin{problemstatement}[Maximum Binary Tree II]
A \textbf{maximum tree} is a tree where every node's value is greater than any other value in its
subtree. You are given the root of a maximum tree built by $\text{Construct}(a)$: the root is
$\max(a)$, its left child is $\text{Construct}(\text{left part})$ and its right child
$\text{Construct}(\text{right part})$. You are not given $a$, only the root, and an integer $\mathit{val}$.
Let $b$ be $a$ with $\mathit{val}$ \emph{appended to the end} (it is guaranteed $b$ has distinct values).
Return $\text{Construct}(b)$. The number of nodes is in $[1,100]$ and all values are $\le 100$.
\end{problemstatement}

\begin{attemptgate}[Attempt it yourself first -- \lv{4}, give yourself 20 minutes]
Recovering $a$ by an inorder walk and rebuilding is allowed but $\bigO{n}$-with-a-full-rebuild and
obscures the structure. The treap view is one recursion.
\peek{$\mathit{val}$ is appended, so it is the \emph{largest index} in $b$. In the value-priority max-treap,
where can a new largest key ever be attached?}
\nudge{Only down the right spine: every existing key is smaller, so the new node descends through right
children until it finds a smaller node -- it then becomes that node's parent, with the old right subtree
as its left child. If it beats the root, it becomes the new root with the whole tree on its left.}
\fullhint{$\text{insert}(t,\mathit{val})$: if $t$ is null or $\mathit{val}>t.\mathit{val}$, return
$\text{node}(\mathit{val},\ \text{left}=t)$; else $t.\text{right}\leftarrow\text{insert}(t.\text{right},
\mathit{val})$ and return $t$. One pass down the right spine, $\bigO{h}$ worst case, $\bigO{n}$ here.}
\end{attemptgate}

\solstep{Understanding the Problem}
A maximum tree is a max-Cartesian tree of $a$: inorder $=$ array order, heap $=$ max on value --
\propref{c26-carters-treap} with the heap flipped and values as keys-priority. Appending $\mathit{val}$
adds the \emph{rightmost} index. The BST-order rule then forces $\mathit{val}$ to be $\ge$ (in index) of
every existing key, so its descent path can only ever take \emph{right} children; and the max-heap rule
makes it stop at the first node whose value is below $\mathit{val}$, becoming that node's parent and
adopting that node's (smaller-valued, larger-indexed) right subtree as its own left child.

\solstep{Key Observations}
\begin{observation}
Because $\mathit{val}$ is a new maximum \emph{index}, it is compared only against right-spine nodes --
never descends left -- so the insertion is a walk down the right spine.
\end{observation}
\begin{observation}
$\mathit{val}$ becomes the parent of the first right-spine node smaller than it (or the root's parent if
it beats the root), and inherits that node's existing right subtree as its left child -- exactly a
treap-priority insertion for the largest key.
\end{observation}
\begin{observation}
This is the recursive one-liner: null-or-$\mathit{val}>$ root $\Rightarrow$ new root with old tree on the
left; else recurse into the right child.
\end{observation}

\solstep{Think It Through With Me}
\begin{thinkbox}
\deadend{Inorder-dump the tree to recover $a$, append $\mathit{val}$, rebuild from scratch.}
{Correct and $\bigO{n}$ (LC 654-style $\text{Construct}$), but it rebuilds nodes that were already right
and misses the point -- appending the largest key only disturbs the right spine. Fine as a fallback, not
as understanding.}
\deadend{Insert $\mathit{val}$ as a plain BST by value, then re-heap.}
{Two trees -- BST order here is on the \emph{index}, not the value; a value-only descent lands $\mathit{val}$
in the wrong place relative to equal-index constraints.}
\aha{Descend the right spine only (the new node is the largest index); stop at the first node smaller than
$\mathit{val}$; make $\mathit{val}$ its parent with the node as $\mathit{val}$'s left child. One recursion,
$\bigO{n}$ worst, and it is a treap insert wearing our \propref{c26-carters-treap} lens.}
\end{thinkbox}

\solstep{Solution}
Define $\text{insert}(t, \mathit{val})$: if $t$ is null or $\mathit{val} > t\to\mathit{val}$, create a node
for $\mathit{val}$ whose left child is $t$ (the entire subtree $t$ lies to the left of the new largest
index and below the new max value) and return it; otherwise set $t\to\mathit{right} =
\text{insert}(t\to\mathit{right}, \mathit{val})$ and return $t$. Call it on the root. Because $\mathit{val}$
is the maximum index, the "$t$ is null" case is only reached at the end of the right spine, and the
"$\mathit{val} > t\to\mathit{val}$" case is the first right-spine node smaller than $\mathit{val}$.

\solstep{Proof of Correctness}
\begin{proofbox}[The right-spine descent yields $\text{Construct}(b)$]
Let $u$ be the first node on the right spine of $\text{Construct}(a)$ whose value is below $\mathit{val}$,
or the null after the spine's last node if all are larger. Every right-spine node above $u$ (closer to the
root) has value $>\mathit{val}$, so by the max-heap rule they stay ancestors of $\mathit{val}$; $u$ and its
whole subtree lie in a range of indices strictly left of $\mathit{val}$'s new index, so $\mathit{val}$, a
larger value and larger index, must become the parent of $u$, keeping $u$ as its left child (the
larger-index side of $\mathit{val}$ is empty since it is last). The recursion returns a new node $v$ with
$v\to\mathit{left}=u$, grafted as $u\to\mathit{parent}\to\mathit{right}$ -- precisely the tree
$\text{Construct}(b)$, whose root/max over $[u\text{ side},\dots,\mathit{val}]$ is $\mathit{val}$. When
$\mathit{val}$ beats the root, $u=$ root and $v$ becomes the new root with the old tree on its left -- again
matching $\text{Construct}(b)$ since $\mathit{val}$ is then $\max(b)$. The algorithm's output is
$\text{Construct}(b)$ in all cases. \end{proofbox}

\solstep{Implementation}
\begin{lstlisting}[caption={Maximum Binary Tree II -- one descent down the right spine (LeetCode 998)}, label={lst:c26-mbt2-soln}]
class Solution {
public:
    TreeNode* insertIntoMaxTree(TreeNode* root, int val) {
        if (!root || root->val < val)          // val is a new max index: only the right spine matters
            return new TreeNode(val, root, nullptr);  // old tree hangs on the new node's LEFT
        root->right = insertIntoMaxTree(root->right, val);  // descend the right spine
        return root;
    }
};
\end{lstlisting}

\solstep{Trace}
\begin{tracebox}[Official sample 3: root \texttt{[5,2,3,null,1]}, \texttt{val = 4}]
$a=[2,1,5,3]$; root $5$ (left $2$, which has right child $1$; right $3$). Call on root $5$: $4>5$? no, so
recurse right on node $3$. At node $3$: $4>3$? yes -- return a new node $4$ with left child $3$. Attach:
$5\to$right $=4$ (with $4\to$left $=3$). Result tree: root $5$, left $2$ (right child $1$), right $4$ (left
child $3$), i.e.\ level order $[5,2,4,\text{null},1,3]$ -- exactly the expected output. $\mathit{val}=4$
descended only through right children, as \propref{c26-carters-treap} guarantees for a largest index. \checkmark
\end{tracebox}

\solstep{Complexity}
\begin{complexitybox}
Time $\bigO{h}$ where $h$ is the tree height: the recursion follows the right spine once. Worst case
$\bigO{n}$ (a right-going chain); on a random maximum tree it is $\bigO{\log n}$ by
\thmref{c26-randbst}/\thmref{c06-exptheight}. Memory $\bigO{h}$ for the call stack (turn it into a loop
over $t\to$right if a deep tree risks RE, exactly the Trap~2 caution from Chapter~7). Verdict: \ac.
\end{complexitybox}

\solstep{Pitfalls and Variants}
\begin{warningbox}
The base case is ``$\mathit{val} > \mathit{root}\to$val'' \emph{or} root null -- both mean "make a new node
here with the current subtree on its left", never on its right (the current subtree is at smaller indices,
and $\mathit{val}$ is the last index). Confusing left/right re-parents the old tree wrongly. \emph{Variant:}
inserting at the \emph{front} would descend the \emph{left} spine by the identical argument; and LC 654
(build a max tree from an array) is the divide-by-maximum construction whose stack build is
\coderef{c26-stackbuild} flipped to a max-heap -- the same tree from the other side.
\end{warningbox}

\begin{insightbox}[Key takeaways]
\begin{enumerate}[nosep, leftmargin=1.6em]
  \item \textbf{Cartesian tree $=$ treap} once you rename value $\to$ priority and index $\to$ key
        (\propref{c26-carters-treap}); Chapter~8's uniqueness gives \thmref{c26-unique} for free.
  \item \textbf{Random permutation $\Rightarrow$ random BST $\Rightarrow$ log height} -- import
        \thmref{c06-exptheight}, never re-prove it (\thmref{c26-randbst}).
  \item \textbf{Appending the max index confines you to one spine} -- that is the whole of LeetCode~998
        (\probref{c26-mbt2}), one recursion deep.
\end{enumerate}
\end{insightbox}

\begin{hardeningbox}
\begin{itemize}[nosep, leftmargin=1.4em]
  \item \textbf{Prefer the loop} -- rewrite the recursion as an explicit walk down $t\to$right to dodge the
        $\bigO{h}$ call stack on a chain, the same hardening as any Chapter-7 BST descent.
  \item \textbf{Front insertion} -- mirror the argument on the left spine; a good drill to be sure you know
        which spine a given key belongs to.
  \item \textbf{Treap by hand} -- build a treap by inserting (key, random priority) pairs and confirm you
        get the Cartesian tree \emph{without a single rotation}: priorities, not rotations, set the shape.
\end{itemize}
\end{hardeningbox}

\begin{recallbox}[10-second refresher -- 9.4 The Cartesian Tree as a Random Treap]
Rename the Cartesian tree's value to a \emph{priority} and its index to a \emph{key} and it is literally a
treap (\propref{c26-carters-treap}); conversely a plain BST built by an insertion order is a Cartesian tree
with priority $=$ arrival time. So a random permutation is a random treap, and Chapter~8's
\thmref{c06-exptheight} hands you expected $\bigO{\log n}$ height for free (\thmref{c26-randbst}). No new
balance proof, ever -- that is the point. The payoff problem, LC~998 (\probref{c26-mbt2}), is a treap
insert of a new largest key down the right spine.
\selftest{
  \item Give the two field-name swaps that turn a Cartesian tree into a treap, and vice-versa.
  \item Why does \propref{c26-carters-treap} re-prove \thmref{c26-unique}?
  \item Why is a random-permutation BST balanced in expectation, and which theorem do you cite?
  \item In LC 998 why can $\mathit{val}$ never descend a left child?
  \item Why is the value-treap of a \emph{sorted} array \emph{not} balanced?
}
\end{recallbox}
